
\section{EEL 5544 : Motivation}
\subsection{What is \emph{random}?}
\begin{itemize}
  \item Some behavior that we cannot perfectly predict, as oppose to
    being \emph{deterministic}.
\end{itemize}

\subsection{Why random?}
\begin{itemize}
\item \textbf{Theologian:}  We can't know God's mind! 
\item \textbf{Copenhagen quantum physicist:} The laws of quantum physics are fundamentally random.
\end{itemize}
  
\subsection{What to do with random?}
Although we cannot perfectly determine a random behavior, we still want to be
able to obtain to some degree its specifications:
\begin{itemize}
  \item \textbf{Probability Theory} is the theory that systematically studies and characterizes randomness.
  \item We will go through an engineering (more operational, less
    about math vigor) treatment of Probability Theory in this course.
\end{itemize}


\subsection{How to study probability?}
\begin{enumerate}
  \item Intuition\\
E.g.) Dr. Wong is a nice gry, the will probably give me a good grade in EEL 5544 .

\begin{itemize}
  \item[-] Not systematic , quessing, cannot quantify.
\end{itemize}

  \item Relative Frequency/Counting
\end{enumerate}

Lig.) Dr. Worg taught EEL5544 3 times in the past There were altogether 200 students who took Pr. War classes. Out of these 200 students, 50 of them got the grode "A". Thas the prob. of getting \(A\) in Dr. Wing's class is \(\frac{50}{200}=\frac{1}{4}\).

\begin{itemize}
  \item[-] Assume equally likely. this may not be the cate.
  \item[-] What about infinite possiblities.
\end{itemize}

\begin{enumerate}
  \setcounter{enumi}{2}
  \item Axiomatic throy
\end{enumerate}

\begin{itemize}
  \item[-] Identify a few basic axious that fit intertion
  \item[-] Build wath ansisteme system based on axioms
  \item[-] Utilize measure thery (a generalization of integration
  \item[-] Developed by Kolnogorou in conly 1900 s
\end{itemize}

